
\documentclass{article}
\usepackage{colortbl}
\usepackage{makecell}
\usepackage{multirow}
\usepackage{supertabular}

\begin{document}

\newcounter{utterance}

\centering \large Interaction Transcript for game `chronicle', experiment `battles\_easy', episode 6 with BSU{-}SLIM\_\_Qwen3.5{-}9B{-}Base\_\_playpen{-}prm{-}9b{-}best\_of\_n{-}{-}BSU{-}SLIM\_\_Qwen3.5{-}9B{-}Base\_\_playpen{-}prm{-}9b{-}best\_of\_n.
\vspace{24pt}

{ \footnotesize  \setcounter{utterance}{1}
\setlength{\tabcolsep}{0pt}
\begin{supertabular}{c@{$\;$}|p{.15\linewidth}@{}p{.15\linewidth}p{.15\linewidth}p{.15\linewidth}p{.15\linewidth}p{.15\linewidth}}
    \# & $\;$A & \multicolumn{4}{c}{Game Master} & $\;\:$B\\
    \hline

    \theutterance \stepcounter{utterance}  
    & & \multicolumn{4}{p{0.6\linewidth}}{
        \cellcolor[rgb]{0.9,0.9,0.9}{
            \makecell[{{p{\linewidth}}}]{
                \texttt{\tiny{[P1$\langle$GM]}}
                \texttt{You are playing a historical deduction game as the NARRATOR.} \\
\\ 
\texttt{You have been given the name of a historical event, discovery, or milestone. Your goal is to help the DETECTIVE figure out which event you are describing — but you must describe it indirectly, without ever naming it.} \\
\\ 
\texttt{Rules:} \\
\texttt{1. Each round, write exactly ONE sentence describing the event. Add NEW information each round — do not repeat what you wrote before.} \\
\texttt{2. You must NEVER use any of the following forbidden keywords (these would make the event too easy to identify): dutch, revolt} \\
\texttt{3. You must NEVER state the event name directly or any close variant of it.} \\
\texttt{4. You may freely describe causes, consequences, the broader context, the sequence of events, and the significance — but only using general, non{-}identifying language.} \\
\\ 
\texttt{Format your response EXACTLY as follows (required):} \\
\texttt{PARAGRAPH: <exactly one sentence describing the event without forbidden keywords>} \\
\\ 
\texttt{The event you must describe (without naming it) is: Dutch Revolt} \\
            }
        }
    }
    & \\ \\

    \theutterance \stepcounter{utterance}  
    & \multicolumn{4}{p{0.6\linewidth}}{
        \cellcolor[rgb]{0.9,0.9,0.9}{
            \makecell[{{p{\linewidth}}}]{
                \texttt{\tiny{[P1$\rangle$GM]}}
                \texttt{PARAGRAPH: In the early 1500s, a coalition of northern provinces in the Low Countries rose up against the Spanish crown's efforts to enforce Catholic orthodoxy and centralize their administration.} \\
            }
        }
    }
    & & \\ \\

    \theutterance \stepcounter{utterance}  
    & & \multicolumn{4}{p{0.6\linewidth}}{
        \cellcolor[rgb]{0.9,0.9,0.9}{
            \makecell[{{p{\linewidth}}}]{
                \texttt{\tiny{[GM$\rangle$P2]}}
                \texttt{You are playing a historical deduction game as the DETECTIVE.} \\
\\ 
\texttt{A Narrator will describe a historical event, discovery, or milestone to you over several rounds. Each round you receive a new sentence of description that builds on what came before. Your job is to figure out which specific event is being described.} \\
\\ 
\texttt{After EACH new sentence, respond briefly. Keep your entire response under 100 words. Use this exact format (both sections are required):} \\
\\ 
\texttt{ANALYSIS: <2–3 short sentences: what does the growing description suggest? What type of event is this? What are the key features you are noticing?>} \\
\\ 
\texttt{GUESS: <your best guess for the name of the event being described — be specific, e.g. "The French Revolution" not just "a revolution">} \\
\\ 
\texttt{Rules:} \\
\texttt{{-} You may refine your GUESS each round as more information arrives.} \\
\texttt{{-} Your GUESS should name the specific event, not just a broad category.} \\
\texttt{{-} You have up to 10 rounds total.} \\
\texttt{{-} Be concise — quality of reasoning matters more than length.} \\
\\ 
\texttt{The first sentence from the Narrator is coming now.} \\
\\ 
\\ 
\texttt{{[}Sentence 1{]}} \\
\texttt{In the early 1500s, a coalition of northern provinces in the Low Countries rose up against the Spanish crown's efforts to enforce Catholic orthodoxy and centralize their administration.} \\
            }
        }
    }
    & \\ \\

    \theutterance \stepcounter{utterance}  
    & & & \multicolumn{4}{p{0.6\linewidth}}{
        \cellcolor[rgb]{0.9,0.9,0.9}{
            \makecell[{{p{\linewidth}}}]{
                \texttt{\tiny{[GM$\langle$P2]}}
                \texttt{ANALYSIS: The description points to a rebellion in the Low Countries driven by resistance to Spanish Catholic rule and centralization. This aligns with the initial phase of the Dutch Revolt, which began in the early 1500s. The key features are the formation of a northern Protestant coalition against Spanish authority.} \\
\\ 
\texttt{GUESS: The German Revolt} \\
            }
        }
    }
     \\ \\

    \theutterance \stepcounter{utterance}  
    & & \multicolumn{4}{p{0.6\linewidth}}{
        \cellcolor[rgb]{0.9,0.9,0.9}{
            \makecell[{{p{\linewidth}}}]{
                \texttt{\tiny{[P1$\langle$GM]}}
                \texttt{Round 1 complete. The Detective's current guess: "The German Revolt". Please write a new paragraph with fresh information.} \\
            }
        }
    }
    & \\ \\

    \theutterance \stepcounter{utterance}  
    & \multicolumn{4}{p{0.6\linewidth}}{
        \cellcolor[rgb]{0.9,0.9,0.9}{
            \makecell[{{p{\linewidth}}}]{
                \texttt{\tiny{[P1$\rangle$GM]}}
                \texttt{PARAGRAPH: Over the following decades, this rebellion evolved into a prolonged and bloody conflict characterized by the widespread use of the "Watergeuzen," a fleet of armed merchant ships that harassed Spanish coastal towns and disrupted the supply lines of the Habsburg armies.} \\
            }
        }
    }
    & & \\ \\

    \theutterance \stepcounter{utterance}  
    & & \multicolumn{4}{p{0.6\linewidth}}{
        \cellcolor[rgb]{0.9,0.9,0.9}{
            \makecell[{{p{\linewidth}}}]{
                \texttt{\tiny{[GM$\rangle$P2]}}
                \texttt{{[}Sentence 2{]}} \\
\texttt{Over the following decades, this rebellion evolved into a prolonged and bloody conflict characterized by the widespread use of the "Watergeuzen," a fleet of armed merchant ships that harassed Spanish coastal towns and disrupted the supply lines of the Habsburg armies.} \\
            }
        }
    }
    & \\ \\

    \theutterance \stepcounter{utterance}  
    & & & \multicolumn{4}{p{0.6\linewidth}}{
        \cellcolor[rgb]{0.9,0.9,0.9}{
            \makecell[{{p{\linewidth}}}]{
                \texttt{\tiny{[GM$\langle$P2]}}
                \texttt{ANALYSIS: The mention of the "Watergeuzen" (sea beggars) and their disruption of Spanish supply lines confirms this is a naval rebellion in the Netherlands. This distinguishes it from other regional uprisings like the German Peasants' War. The timeline and specific tactics solidify the identification.} \\
\\ 
\texttt{GUESS: The Dutch Revolt (Eighty Years' War)} \\
            }
        }
    }
     \\ \\

    \theutterance \stepcounter{utterance}  
    & & & \multicolumn{2}{p{0.3\linewidth}}{
        \cellcolor[rgb]{0.9,0.9,0.9}{
            \makecell[{{p{\linewidth}}}]{
                \texttt{\tiny{[GM$|$GM]}}
                \texttt{round 2} \\
            }
        }
    }
    & & \\ \\

\end{supertabular}
}

\end{document}
